% SPDX-FileCopyrightText: Copyright (c) 2016-2026 Objectionary.com
% SPDX-License-Identifier: MIT

\begin{proof}[Proof of \cref{thm:contextualization-total}]
Fix an absolute \(k\) and argue by structural induction on the expression \(e\),
  establishing at once that a derivation of \(\phinoContextualize{e}{k}{e_1}\)
  exists for some \(e_1\) and that it is finite.
By the grammar of \cref{fig:ebnf} an expression is a locator ($Q$ or
  \(\phiTerminal{\xi}\)), a terminator $T$, a formation, a dispatch, or an
  application, and its outermost constructor tells these cases apart.
The rules of \cref{fig:contextualization} match the cases one-to-one, so exactly
  one rule applies to every expression.
For $Q$, \(\phiTerminal{\xi}\), $T$ and a formation $[[ B ]]$
  the matching rule---\nameref{r:cg}, \nameref{r:cxi}, \nameref{r:ct} or
  \nameref{r:cf}---is an axiom that yields a result in a single step, so the base
  cases hold; \nameref{r:cf} rewrites the formation to itself and never descends
  into \(B\), halting the recursion at every formation boundary.
A dispatch $e_1.\tau$ selects \nameref{r:cd}, whose single premise
  contextualizes the subject \(e_1\); an application $e_1( \tau -> e_2 )$ or
  $e_1( \phiTerminal{\alpha_i} -> e_2 )$ selects \nameref{r:ca} or
  \nameref{r:caa}, whose two premises contextualize \(e_1\) and \(e_2\).
Every premise is a strict subexpression of its conclusion, so by the induction
  hypothesis each has a finite derivation and a result, and the chosen rule
  assembles them into a finite derivation of the whole; the argument \(k\) is
  copied wholesale at \nameref{r:cxi} and never traversed.
The recursion therefore follows the finite syntax tree and every rule is defined
  on its case, so contextualization is total on \(\mathcal{E} \times \mathcal{K}\)
  and each derivation terminates.
\end{proof}
